\subsection{"壮阳"补充剂与性健康产品：循证评价}

市面上"壮阳""补肾""延时"的保健品与草本补充剂种类繁多，宣传往往远大于证据。本节按\textbf{证据强度}分类说明，帮助读者少花冤枉钱、少担健康风险。

\vspace{0.5em}
\noindent\textbf{1. 证据较弱或矛盾的常用成分}

\begin{itemize}
	\item \textbf{玛卡（Maca）}：部分小型研究显示可能改善性欲，但设计质量参差、结论不一致，尚无高质量证据支持其治疗性功能障碍。作为食品补充通常安全。
	\item \textbf{育亨宾（Yohimbine）}：来自一种树皮，历史上用于勃起功能障碍。效果证据有限，且可引起\textbf{焦虑、心悸、血压升高}，与抗抑郁药等联用有风险，\textbf{不推荐自行服用}。
	\item \textbf{L-精氨酸 / L-瓜氨酸}：作为一氧化氮前体，理论上可能轻度改善勃起，小型研究结果不一；对健康人价值有限，腹胀腹泻常见。
	\item \textbf{锌、维生素 D、复合维生素}：仅在\textbf{确有缺乏}时才可能有益；\textbf{不缺的人额外补充通常不提升性功能}，过量补锌反而有害。
	\item \textbf{DHEA（脱氢表雄酮）}：证据不足以支持常规用于提升性欲；可能影响激素水平，需医生评估。
	\item \textbf{人参、鹿茸、淫羊藿等传统药材}：传统上归于"壮阳"，现代研究证据总体薄弱或矛盾，个体差异大，且可能与其他药物相互作用。
\end{itemize}

\noindent\textbf{2. 真正的风险：违规添加}

\begin{itemize}
	\item 不少网售"壮阳保健品""延时喷剂"被检出\textbf{违规添加西地那非、他达拉非或其类似物}，剂量不可控、来源不明。
	\item 正在服用\textbf{硝酸酯类药物}（如硝酸甘油）者误服这些产品，可能引发\textbf{严重低血压甚至危及生命}。
	\item 这类产品还可能\textbf{掩盖真实病情}，延误对勃起功能障碍、心血管疾病的规范诊治。
\end{itemize}

\noindent\textbf{3. 正确的态度}

\begin{itemize}
	\item \textbf{食物不能替代药物}：均衡饮食有益整体健康，但不存在"特效壮阳食物"（生蚝、韭菜、腰子）。
	\item \textbf{先找原因，再谈补充}：勃起困难、性欲低下常与血管、神经、内分泌、睡眠、情绪或药物有关，应先由医生评估，而非先用保健品试错。
	\item \textbf{如需用药，用"药"而非"保健品"}：确有适应证时，PDE5 抑制剂等处方药在医生指导下使用，安全性与有效性都有保障。
	\item \textbf{就诊前列清单}：告诉医生你正在服用的所有补充剂；正在服用硝酸酯类、$\alpha$ 受体阻滞剂者尤其要谨慎。
\end{itemize}

\begin{tcolorbox}[colback=yellow!10!white,colframe=orange!70!black,title=科学提示]
判断一款"壮阳产品"是否可信，最快的办法是看它是否宣称"纯天然、无副作用、立刻见效、不用看医生"。\textbf{真正有效的治疗，从不回避"请就医"这三个字}。与其在保健品上反复试错，不如把同样的钱和精力用在睡眠、运动、减压和一次规范的医学评估上。
\end{tcolorbox}
